\chapter{External-input provenance audit}
\label{chap:provenance-audit}

This audit chapter records source identifiers, locators, inventories, claim ledgers, and the trust boundary.  It is technical bookkeeping, not a mathematical prerequisite and not a parent of every chapter.

\section{Source-bearing input structures}

% SOURCE: PROJECT_BOUNDARY.md, External inputs; Acceptance criteria.
% VERIFIED-IN: KIP126/External/Provenance.lean, declaration SourceId.
\begin{definition}[Source identifier]
  \label{def:source-id}
  \lean{KIP126.External.SourceId}
  \leanok
  A source identifier is a finite, stable key for an item in the project
  inventory; the identifier itself carries no mathematical claim.  Its
  associated source-kind entry records whether the item is a paper, prior
  literature result, machine archive, or governance record.
  Human-readable citation data and artifact hashes live in the external
  inventory rather than in this index type.
\end{definition}

% SOURCE: PROJECT_BOUNDARY.md, External inputs; Acceptance criteria.
% VERIFIED-IN: KIP126/External/Provenance.lean, declaration SourceKind.
\begin{definition}[Source kind]
  \label{def:source-kind}
  \lean{KIP126.External.SourceKind}
  \leanok
  A source kind records whether an inventory entry is a paper, literature
  source, machine computation, or project-governance record.
\end{definition}

% SOURCE: PROJECT_BOUNDARY.md, External inputs; Acceptance criteria.
% VERIFIED-IN: KIP126/External/Provenance.lean, declaration Locator.
\begin{definition}[Source locator]
  \label{def:locator}
  \lean{KIP126.External.Locator}
  \leanok
  A source locator is a structured description string and may optionally
  name the artifact in which the description is found.  The locator records
  where to look; it does not by itself assert that the located statement is
  true.  The description is explanatory ledger/human-audit metadata: when
  available, the surrounding claim prose interprets it as a theorem, page,
  table-row, file-range, or output key.  The machine checker validates the
  associated path, not the prose description itself.
\end{definition}

% SOURCE: PROJECT_BOUNDARY.md, External inputs; Acceptance criteria.
% VERIFIED-IN: KIP126/External/Provenance.lean, declaration SourceRef.
\begin{definition}[Source reference]
  \label{def:source-ref}
  \lean{KIP126.External.SourceRef}
  \leanok
  A source reference pairs a stable source identifier with a structured locator and
  may carry an optional note.  It identifies the provenance of one supplied
  proposition.  Citation metadata, source class, and content hashes are
  resolved by the inventory keyed by the source identifier; they are not
  silently inferred from the reference.
  \uses{def:source-id,def:locator}
\end{definition}

% SOURCE: PROJECT_BOUNDARY.md, External inputs; Acceptance criteria.
% VERIFIED-IN: KIP126/External/Provenance.lean, declaration ArtifactRef.
\begin{definition}[Artifact reference]
  \label{def:artifact-ref}
  \lean{KIP126.External.ArtifactRef}
  \leanok
  An artifact reference records a local or archived path, a SHA-256 digest,
  and an optional version.  It is metadata for a reproducibility check and is
  not itself evidence for a mathematical proposition.
\end{definition}

% SOURCE: PROJECT_BOUNDARY.md, External inputs and Acceptance criteria.
% VERIFIED-IN: KIP126/External/Provenance.lean, declaration EvidenceMetadata.
\begin{definition}[Evidence metadata]
  \label{def:evidence-metadata}
  \lean{KIP126.External.EvidenceMetadata}
  \leanok
  Evidence metadata consists of the method string used to obtain or transcribe
  a finite input and optional artifact metadata.  It is a reusable view of
  metadata; the proposition supported by that metadata remains a field of an
  explicit \texttt{ExternalEvidence} value.
  \uses{def:artifact-ref}
\end{definition}

% SOURCE: PROJECT_BOUNDARY.md:149--167, explicit proposition-bearing wrappers.
% VERIFIED-IN: KIP126/External/Provenance.lean, declaration ExternalResult.
\begin{definition}[External result]
  \label{def:external-result}
  \lean{KIP126.External.ExternalResult}
  \leanok
  For a proposition $P$, an external result is a pair $(p,\sigma)$ with
  $p:P$ and a source reference $\sigma$.  It is supplied as an explicit
  argument to every theorem that uses it; the project shall not install it as
  an axiom or an untracked global instance.
  \uses{def:source-ref}
\end{definition}

% SOURCE: PROJECT_BOUNDARY.md:149--167, explicit proposition-bearing wrappers.
% VERIFIED-IN: KIP126/External/Provenance.lean, declaration ExternalEvidence.
\begin{definition}[External evidence]
  \label{def:external-evidence}
  \lean{KIP126.External.ExternalEvidence}
  \leanok
  For a proposition $P$, external evidence is a record containing evidence
  $p:P$, a source reference, a method string, and optional artifact metadata.
  It is used for finite Ext data, differentials, table entries, and diagram
  relations.  Like an external result, it is always an explicit theorem
  parameter.  The wrapper records the declared evidence and its provenance;
  it does not turn either one into a project theorem.
  \uses{def:source-ref,def:artifact-ref}
\end{definition}

\section{Inventory validity and trust boundary}

% VERIFIED-IN: KIP126/External/SourceInventory.lean, declaration SourceAvailability.
\begin{definition}[Source availability]
  \label{def:source-availability}
  \lean{KIP126.External.SourceAvailability}
  \leanok
  A source-availability value records four Boolean flags for metadata, PDF,
  extracted-text, and source-archive availability.  These fields are claims
  about acquisition state, not mathematical evidence.  For canonical rows,
  the filesystem checker compares them with the status file and the artifacts
  actually present; the typed inventory separately proves that each row's
  status class agrees with its flags.
\end{definition}

% VERIFIED-IN: KIP126/External/SourceInventory.lean, declaration SourceStatusClass.
\begin{definition}[Source status class]
  \label{def:source-status-class}
  \lean{KIP126.External.SourceStatusClass}
  \leanok
  The status class is the stable acquisition summary shared with the JSON
  ledger: source of record, full text, metadata only, or partial availability.
  Its encoder and decoder are inverse on the finite set of admitted classes.
\end{definition}

% VERIFIED-IN: KIP126/External/SourceInventory.lean, declaration SourceEntry.
\begin{definition}[Typed source entry]
  \label{def:source-entry}
  \lean{KIP126.External.SourceEntry}
  \leanok
  A typed source entry exposes the JSON/Lean synchronization fields: source
  identifier, kind, citation keys, local directory, status-file path, status
  class, and availability flags.  Its validity predicate checks schema-safe,
  nonempty, duplicate-free citation keys; a nonempty safe
  directory; kind/status consistency; the distinguished source-of-record
  convention; and a safe, source-relative status path when present.
  Filesystem existence, bibliographic prose, and artifact hashes remain under
  the external checker.
  \uses{def:source-id,def:source-kind,def:source-status-class,def:source-availability}
\end{definition}

% SOURCE: PROJECT_BOUNDARY.md, External inputs and Acceptance criteria.
% VERIFIED-IN: KIP126/External/SourceInventory.lean, declaration SourceInventory.inventory.
\begin{definition}[Source inventory]
  \label{def:source-inventory}
  \lean{KIP126.External.SourceInventory.inventory}
  \leanok
  The source inventory is a complete typed projection keyed by
  Definition~\ref{def:source-id}.  Its Lean rows expose source kinds,
  citation keys, local status-file paths, and acquisition-status classes.  The
  machine-readable ledger additionally supplies titles, canonical identifiers,
  artifact paths, hashes, and filesystem verification; a separate
  checker validates those fields.  Neither representation proves the
  propositions carried by an \texttt{ExternalResult} or
  \texttt{ExternalEvidence} value.  The Lean resolver uses a
  \texttt{SourceRef}'s stable identifier to select its catalogue row.
  \uses{def:source-entry}
\end{definition}

% VERIFIED-IN: KIP126/External/SourceInventory.lean, declaration SourceRef.InventoryValid.
\begin{definition}[Inventory-valid source reference]
  \label{def:inventory-valid-source-ref}
  \lean{KIP126.External.SourceRef.InventoryValid}
  \leanok
  A source reference is inventory-valid when it is structurally valid and any
  artifact named by its locator is a syntactically safe relative path lying
  below the directory of the selected source row (a syntactic directory
  prefix; filesystem links are outside Lean's view).  Thus absolute paths,
  backslashes, empty components, and the components \texttt{.} and
  \texttt{..} are rejected.  This predicate does not inspect the filesystem
  and does not claim that the path is an artifact listed in the JSON ledger;
  existence, canonical-ledger membership, and digest equality remain
  obligations of the external checker.
  \uses{def:source-ref,def:source-inventory}
\end{definition}

% VERIFIED-IN: KIP126/External/SourceInventory.lean, declaration ExternalEvidence.InventoryValid.
\begin{definition}[Inventory-valid external evidence]
  \label{def:inventory-valid-evidence}
  \lean{KIP126.External.ExternalEvidence.InventoryValid}
  \leanok
  An evidence wrapper is inventory-valid when its structural metadata is
  valid, its source locator is safe and source-relative, and an optional
  attached artifact has a lower-case 64-hex SHA-256 spelling and a safe path
  below the same source directory.  This Lean predicate intentionally stops at
  syntax and syntactic source-relative containment.  For canonical claim
  locators, the JSON checker is responsible for artifact-list membership, file
  existence, and equality with the checked-in digest.  An arbitrary dynamic
  wrapper is not visible to that checker, and its own digest text is not
  automatically compared unless the caller exports and checks it.
  \uses{def:external-evidence,def:inventory-valid-source-ref,def:artifact-ref}
\end{definition}

% VERIFIED-IN: KIP126/External/SourceInventory.lean, declaration SourceInventory.inventory_complete.
\begin{theorem}[Source inventory completeness]
  \label{thm:source-inventory-complete}
  \lean{KIP126.External.SourceInventory.inventory_complete}
  \leanok
  The ordered Lean rows project to exactly the closed list of source
  identifiers, so no admitted source is silently omitted.
  \uses{def:source-inventory,def:source-id}
\end{theorem}

\begin{proof}
  Map the canonical lookup rows over the finite \texttt{SourceId.all} list.
  The Lean equality \texttt{SourceInventory.inventory\_complete} is the
  resulting catalogue invariant; duplicate-freeness is recorded separately by
  \texttt{SourceInventory.inventory\_complete\_nodup}.
  \leanok
\end{proof}

% VERIFIED-IN: KIP126/External/SourceInventory.lean, declaration SourceInventory.inventory_globally_valid.
\begin{theorem}[Global source-inventory validity]
  \label{thm:source-inventory-globally-valid}
  \lean{KIP126.External.SourceInventory.inventory_globally_valid}
  \leanok
  Across all rows of the canonical source inventory, source directories are
  pairwise distinct, present status-file paths are pairwise distinct, and all
  citation keys are pairwise distinct.
  \uses{def:source-inventory}
\end{theorem}

\begin{proof}
  The three flattened finite lists are computed from the 18 canonical rows;
  their no-duplicate predicates reduce to decidable equalities of the checked
  strings.
  \leanok
\end{proof}

% VERIFIED-IN: KIP126/External/SourceInventory.lean, declaration SourceInventory.inventory_valid.
\begin{theorem}[Source inventory row validity]
  \label{thm:source-inventory-valid}
  \lean{KIP126.External.SourceInventory.inventory_valid}
  \leanok
  Every row selected by the canonical inventory satisfies the typed kind,
  status, citation-key, and source-relative path invariants.
  \uses{def:source-inventory,def:source-entry}
\end{theorem}

\begin{proof}
  Each of the 18 lookup cases reduces the typed row predicates to decidable
  string, kind, status, key, and directory-prefix checks; the checked
  \texttt{inventory} definition discharges all of them.
  \leanok
\end{proof}

% VERIFIED-IN: KIP126/External/SourceInventory.lean, declaration SourceInventory.projectionManifest.
\begin{definition}[Source projection manifest]
  \label{def:source-projection}
  \lean{KIP126.External.SourceInventory.projectionManifest}
  \leanok
  The projection manifest is the ordered, delimiter-stable export of the
  fields shared by Lean and JSON: source ID, directory, kind, citation keys,
  status-file path, status class, and the four availability flags.  The
  inventory checker executes the Lean exporter and rejects any disagreement
  with those JSON fields.  Titles, publication data, artifact membership,
  filesystem existence, and digests remain separate checker obligations.
  \uses{def:source-inventory,def:source-entry}
\end{definition}

% VERIFIED-IN: KIP126/External/Evidence.lean, declaration ExternalEvidence.withArtifact_inventoryValid.
\begin{theorem}[Inventory-valid artifact transport]
  \label{thm:inventory-valid-artifact-transport}
  \lean{KIP126.External.ExternalEvidence.withArtifact_inventoryValid}
  \leanok
  Replacing an evidence wrapper's optional artifact preserves
  \texttt{InventoryValid} when the original certificate, the new artifact's
  syntactic SHA-256 certificate, and its source-relative path certificate are
  supplied explicitly.  This is a transport lemma, not a filesystem theorem.
  \uses{def:inventory-valid-evidence,def:artifact-ref}
\end{theorem}

\begin{proof}
  Unfold the wrapper update.  The proposition, source locator, and method
  remain unchanged; combine the supplied digest and path certificates with the
  structural validity fields.
  \leanok
\end{proof}

% SOURCE: PROJECT_BOUNDARY.md, Examples, remarks, and questions.
\begin{definition}[Open proposition]
  \label{def:open-proposition}
  \notready
  An open proposition is a proposition recorded as a formalization target but
  neither assumed nor asserted.  Its Blueprint node has no incoming edge to a
  proved conclusion unless a later theorem explicitly takes a proof of it as
  a hypothesis.
\end{definition}

% SOURCE: PROJECT_BOUNDARY.md, Confirmed design decisions and Axiom policy.
\begin{definition}[KIP126 trust boundary]
  \label{def:kip126-trust-boundary}
  \uses{def:external-result,def:external-evidence,def:open-proposition,def:kervaire-main-input}
  \notready
  The KIP126 development may use Lean's foundational axioms and declarations
  proved in the pinned Mathlib revision.  Paper-internal arguments are project
  proof obligations.  Results from earlier literature are explicit external
  results, while high-stem computations and every appendix row are explicit
  external evidence.  The conditional permanent-cycle and geometric endpoints
  are universally quantified over the explicit main-input record of
  Definition~\ref{def:kervaire-main-input}; none of its fields is installed as
  a global instance.  No project declaration may use \texttt{axiom},
  \texttt{sorry}, or \texttt{admit}.
\end{definition}

\section{Claim ledger and catalogued inputs}

% VERIFIED-IN: KIP126/External/Claims.lean, declaration ExternalRootId.
\begin{definition}[External root identifier]
  \label{def:external-root-id}
  \lean{KIP126.External.ExternalRootId}
  \leanok
  An external-root identifier is a closed, enumerable key for a literature
  result, machine catalogue, appendix family, or finite evidence family that
  may cross the project boundary.  It carries no proposition by itself.
\end{definition}

% VERIFIED-IN: KIP126/External/Claims.lean, declaration ExternalClaimClass.
\begin{definition}[External claim class]
  \label{def:external-claim-class}
  \lean{KIP126.External.ExternalClaimClass}
  \leanok
  The five classes are literature result, machine evidence, table evidence,
  transcribed evidence, and composite result.  The result wrapper supports
  exactly literature and composite rows; the evidence wrapper supports exactly
  machine, table, and transcribed rows.  These compatibility predicates
  constrain which wrapper interface may bind to a canonical row; the class
  itself is provenance metadata, not a proof of the associated proposition.
\end{definition}

% VERIFIED-IN: KIP126/External/Claims.lean, declaration ExternalClaimRecord.
\begin{definition}[External claim record]
  \label{def:external-claim-record}
  \lean{KIP126.External.ExternalClaimRecord}
  \leanok
  A claim record binds one root identifier to a trust classification, an owner
  declaration name, a Blueprint-label or enumerated source target, explicit
  dependencies, and a valid \texttt{SourceRef}.  A valid row also has a
  project-namespace owner, a
  recognized target class, no repeated or direct self-dependencies, and an
  inventory-valid locator.  Composite rows retain a nonempty dependency list
  and are not reclassified as primitive external facts.
  \uses{def:external-root-id,def:external-claim-class,def:inventory-valid-source-ref}
\end{definition}

% VERIFIED-IN: KIP126/External/Claims.lean, declaration externalClaimLedger.
\begin{definition}[External claim ledger]
  \label{def:external-claim-ledger}
  \lean{KIP126.External.externalClaimLedger}
  \leanok
  The claim ledger is the checked-in finite record supplying one valid row for
  each enumerated external root.  Its completeness and source-coverage facts
  are stated separately below; this definition is a provenance/catalogue
  object, not a theorem asserting any recorded mathematical claim.
  \uses{def:external-claim-record}
\end{definition}

% VERIFIED-IN: KIP126/External/Claims.lean, declaration externalClaimLedger_complete.
\begin{theorem}[External claim ledger completeness]
  \label{thm:external-claim-ledger-complete}
  \lean{KIP126.External.externalClaimLedger_complete}
  \leanok
  The IDs in the ledger entries enumerate exactly the closed list of external
  roots.  This is a finite catalogue invariant, not a proof of the propositions
  named by those roots.
  \uses{def:external-claim-ledger,def:external-root-id}
\end{theorem}

\begin{proof}
  Case analysis over the closed \texttt{ExternalRootId} index verifies the lookup ID
  equation; mapping the lookup over \texttt{ExternalRootId.all} gives the stated
  equality.  Duplicate-freeness is recorded separately by
  \texttt{externalClaimLedger\_nodup}.
  \leanok
\end{proof}

% VERIFIED-IN: KIP126/External/Claims.lean, declaration externalClaimLedger_refs_nodup.
\begin{theorem}[Unique canonical claim references]
  \label{thm:external-claim-refs-unique}
  \lean{KIP126.External.externalClaimLedger_refs_nodup}
  \leanok
  The exact source references attached to the 55 canonical claim roots are
  pairwise distinct.  Therefore a wrapper whose reference is equal to a
  canonical row cannot ambiguously select two roots.
  \uses{def:external-claim-ledger}
\end{theorem}

\begin{proof}
  The finite list of claim records is mapped to its references and the
  resulting equality checks reduce to decidable equality of source IDs,
  descriptions, and optional artifact paths.
  \leanok
\end{proof}

% VERIFIED-IN: KIP126/External/Claims.lean, declaration externalClaimLedger_no_dependency_cycle.
\begin{theorem}[Acyclic external-claim dependencies]
  \label{thm:external-claim-dependency-acyclic}
  \lean{KIP126.External.externalClaimLedger_no_dependency_cycle}
  \leanok
  The dependency relation of the canonical external-claim ledger admits a
  natural-number rank that strictly decreases along every direct dependency.
  Consequently, the transitive closure of the direct-dependency relation has
  no path from a root back to itself.
  \uses{def:external-claim-ledger}
\end{theorem}

\begin{proof}
  Assign the explicit finite rank \texttt{ExternalRootId.dependencyRank} to
  each of the 55 roots.  Exhaustive reduction of every admitted edge proves
  that a dependency has smaller rank than its owner.  Strict descent in the
  natural numbers excludes cycles.
  \leanok
\end{proof}

% VERIFIED-IN: KIP126/External/Claims.lean, declaration CataloguedExternalResult.
\begin{definition}[Catalogued external result]
  \label{def:catalogued-external-result}
  \lean{KIP126.External.CataloguedExternalResult}
  \leanok
  A catalogued external result consists of an explicit
  \texttt{ExternalResult}, an external-root identifier, equality of the
  result's source reference with that root's canonical claim row, and a proof
  that the row's trust class supports literature-style results.  The canonical
  row validity theorem then supplies the stronger inventory-validity
  certificate.  The proposition and its proof still come from the explicit
  wrapper value, not from the ledger metadata.
  The root, locator, and trust-class fields are metadata checks; they do not
  type-check that the proposition $P$ is definitionally the proposition named
  by the future owner declaration.  Exact source references are pairwise
  distinct across the canonical ledger, so the reference equality does select
  a unique canonical root.
  \uses{def:external-result,def:external-claim-ledger,def:external-claim-class,def:inventory-valid-source-ref,thm:external-claim-refs-unique}
\end{definition}

% VERIFIED-IN: KIP126/External/Claims.lean, declaration CataloguedExternalEvidence.
\begin{definition}[Catalogued external evidence]
  \label{def:catalogued-external-evidence}
  \lean{KIP126.External.CataloguedExternalEvidence}
  \leanok
  A catalogued external-evidence value similarly binds an explicit
  \texttt{ExternalEvidence} to one canonical root and compatible evidence
  class, and carries the stronger inventory-validity certificate for both its
  locator and optional evidence artifact.  If an artifact is attached, its
  path must equal the canonical claim locator.  The Lean certificate checks
  the artifact's safe source-relative path and SHA-256 spelling.  Independently
  of whether a wrapper artifact is attached, the JSON checker validates the
  canonical locator's ledger membership, required status, filesystem
  existence, and checked-in digest; it does not compare an arbitrary wrapper's
  digest text with that JSON value.  The wrapper does not type-check that its
  proposition is definitionally the one named by the future owner declaration.
  \uses{def:external-evidence,def:external-claim-ledger,def:external-claim-class,def:inventory-valid-evidence}
\end{definition}

% VERIFIED-IN: KIP126/External/Claims.lean, declaration externalClaimProjectionManifest.
\begin{definition}[External-claim projection manifest]
  \label{def:external-claim-projection}
  \lean{KIP126.External.externalClaimProjectionManifest}
  \leanok
  The claim projection exports, for every canonical root, its identifier,
  source, optional locator artifact, locator description, owner, and
  Blueprint-label or enumerated source target.  The external checker rejects
  duplicate or missing root rows, duplicate owners or targets, missing source
  coverage, unknown Blueprint/source targets, unlocated claims without the
  required exact fallback locator, and any nonempty canonical locator artifact
  path that does not name a listed, required, existing regular file under the
  corresponding source directory.  This exported metadata still does not
  assert the truth of a claim.
  \uses{def:external-claim-ledger}
\end{definition}

% VERIFIED-IN: KIP126/External/Claims.lean, declaration every_source_has_valid_claim.
\begin{theorem}[Source coverage by claim roots]
  \label{thm:source-coverage-by-claims}
  \lean{KIP126.External.every_source_has_valid_claim}
  \leanok
  Every source identifier in the closed source inventory is referenced by at
  least one valid claim row.  This proves valid coverage of the provenance
  index only; it does not assert the truth of the referenced literature or
  computation.
  \uses{def:external-claim-ledger,def:source-id}
\end{theorem}

\begin{proof}
  For each constructor of \texttt{SourceId}, select the corresponding named claim
  row in the finite ledger.  The resulting 18 witnesses are checked by
  definitional reduction; the canonical ledger then supplies validity of the
  selected row.
  \leanok
\end{proof}
